\documentclass[11pt]{article}
\usepackage{amsmath}
\usepackage{amssymb}
\usepackage{amsfonts}
\usepackage{graphicx}
\usepackage{booktabs}
\usepackage{siunitx}
\usepackage{geometry}
\usepackage{caption}
\usepackage{subcaption}
\usepackage{hyperref}
<<<<<<< HEAD
\geometry{margin=1in}

\title{\textbf{Paper 1: Binding Mechanisms and Cage Stability\\
in the 600-Cell Lattice}\\[6pt]
{\large Standard Model Emergence in the 600-Cell Lattice Series}\\[4pt]
{\normalsize Version 6 --- revised for consistency with the CPP paper series}}

\author{Thomas Lee Abshier, ND \and Grok (xAI) \and Claude Sonnet (Anthropic)\\[4pt]
Hyperphysics Institute\\
\url{https://hyperphysics.com}\\
\texttt{drthomas007@protonmail.com}}

\date{Version 6, 26 March 2026}
=======
\hypersetup{colorlinks=true,linkcolor=blue,citecolor=blue,urlcolor=blue}
\geometry{margin=1in}

\title{\textbf{SM-1: Binding Mechanisms and Cage Stability\\
in the 600-Cell Lattice}\\[6pt]
{\large 600-Cell Standard Model Emergence Series}\\[4pt]
{\normalsize Version 6, 26 March 2026}}

\author{Thomas Lee Abshier, ND \and Grok (xAI) \and Claude Sonnet (Anthropic)}

\date{\normalsize Hyperphysics Institute\\
\url{https://hyperphysics.com}\\
\texttt{drthomas007@protonmail.com}}
>>>>>>> c085da0cc636ab5fe5c35ba45eac2ff84c41fcdc

\begin{document}

\maketitle

\begin{abstract}
In Conscious Point Physics (CPP), Standard Model particles emerge as stable
aggregates of Charged Conscious Points (CPs) organised into geometric
``cages'' within the 600-cell lattice.  This paper derives the binding
mechanisms: Space Stress Vector (SSV) gradients between CPs of opposite
polarity create attractive forces that minimise total SSV energy in
specific symmetric configurations.  We present the SSV force law,
<<<<<<< HEAD
demonstrate energy minimisation for the principal cage geometries,
=======
demonstrate energy minimisation for the principal cage geometries
(tetrahedral, icosahedral, dodecahedral, and 30-vertex shell),
>>>>>>> c085da0cc636ab5fe5c35ba45eac2ff84c41fcdc
and provide a worked numerical example for the electron.

\textbf{Corrections in Version~6} (for consistency with subsequent CPP papers):
The fourth cage for the top quark was previously listed as the C$_{60}$
fullerene (60 vertices); exact 600-cell shell computation (PS-1, March 2026)
shows no 60-vertex distance shell exists.  The leading candidate is the
\textbf{30-vertex shell} at $d^2 = 2$ (shell~3, all vertices degree~4,
vertex-transitive).  The binding energy table and Appendix~B.1 are
updated accordingly.  One calibration constant ($\text{SSV}_0 = 0.2555$~MeV,
fixed to the electron mass) is used; this is noted explicitly throughout.
\end{abstract}

\section*{Keywords}
Conscious Point Physics, 600-cell lattice, binding mechanisms, cage stability,
Space Stress Vector, Zitterbewegung, Dipole Sea, geometric suppression,
Standard Model emergence
<<<<<<< HEAD

=======

%=====================================================================
\section*{Plain Language Summary}
%=====================================================================

\textit{What holds a particle together, and why does it have mass?}

The Standard Model tells us that particles exist and gives precise rules
for how they interact, but it does not explain why particles have the
shapes, charges, or masses they do.  Conscious Point Physics (CPP)
proposes a geometric answer: every particle is a stable cluster of
fundamental charge-carrying points arranged in a specific cage, and the
properties of the particle follow from the geometry of that cage.

The cage sits inside a four-dimensional crystal structure called the
600-cell --- a polytope with 120 vertices connected by edges in a
pattern governed by the golden ratio.  Particles do not move through
this crystal the way a ball rolls across a floor; they are
self-sustaining patterns of organisation within it.

A deeper question underlies all of this: what actually causes a
Conscious Point to move?  Standard physics describes forces with great
precision but treats them as primitive --- a field exists, a charge is
in the field, and motion follows, without any account of the mechanism
behind that motion.  CPP proposes a specific answer: each CP perceives
the signals of nearby CPs, computes the net direction of attraction or
repulsion, and moves a lawfully determined distance each absolute clock
tick.  Motion is not produced by a force acting on a passive object; it
is the result of perception and response by an entity that is, in a
minimal sense, aware of its neighbours.  The Space Stress Vector (SSV)
is the name CPP gives to the aggregate signal that a CP reads from its
local environment.  The full development of this picture --- what CPs
are, how they perceive, what laws govern their response --- is the
subject of the foundational CPP postulates rather than this paper.
What this paper shows is that if you accept those postulates, the
specific cage geometries of the 600-cell lattice produce the binding
energies and stability conditions that correspond to the known particles
of the Standard Model.

What holds a cage together emerges from the SSV field.  Each CP computes
the vector sum of attraction and repulsion signals from all nearby CPs,
and moves accordingly each clock tick.  To an outside observer this
appears as motion along the steepest SSV gradient --- but the gradient
is the aggregate description of many individual displacement decisions,
not the cause of them.  When several CPs settle into a symmetric
arrangement that minimises the total SSV energy through this process,
the result is a stable, self-sustaining cage.

A particle is not simply any region of the lattice that is more
organised than its surroundings.  Kinetic energy, magnetic fields, and
atomic bonds all impose local order on the lattice without constituting
particles.  What distinguishes a particle is that its internal
organisation is self-sustaining: the cage geometry closes on itself with
no loose ends, the Zitterbewegung resonance inside it maintains the
pattern without external support, and the whole structure carries a
definite, quantised rest mass.  For fermions, an unpaired central charge
nucleates this closed pattern --- the asymmetry is what pulls the Dipole
Sea into a stable cage around it.  For bosons and neutrinos, no unpaired
charge is needed; specific configurations of paired charges achieve
self-sustaining resonance through their own symmetry.

The four stable cages correspond to four families of particles.  The
smallest --- four vertices in a tetrahedral arrangement --- is the cage
of the electron and the light quarks.  The next --- twelve vertices in
an icosahedral arrangement --- corresponds to heavier particles such as
the charm quark and tau lepton.  The third --- twenty vertices in a
dodecahedral arrangement --- is the bottom quark and the Higgs boson.
The fourth candidate --- thirty vertices at equal distance from a
central point --- is the top quark, the heaviest known particle.  The
more vertices a cage has, the more SSV energy it stores, and the more
massive the particle.

One thing this paper does \textit{not} do is derive the absolute mass
scale from scratch.  A single calibration constant
($\text{SSV}_0 = 0.2555$~MeV) is fixed to the measured electron mass;
the \textit{hierarchy} among particles comes entirely from cage geometry,
but the overall scale is anchored to experiment.  This is an honest
limitation, clearly noted throughout.

This paper is the foundation for the rest of the series.  SM-2 uses the
cage picture to estimate masses for all Standard Model particles.  SS-1
and SM-3 later provide rigorous proofs --- derived without free
parameters --- of properties that this paper can only motivate
geometrically.

>>>>>>> c085da0cc636ab5fe5c35ba45eac2ff84c41fcdc
%=====================================================================
\section{Introduction}
%=====================================================================

The Standard Model describes particles and their interactions accurately
but provides no geometric origin for their existence, masses, or
symmetries~\cite{pdg2024}.  In Conscious Point Physics (CPP), the universe
is mediated by a finite 4D 600-cell lattice whose 120 vertices (120CPs)
serve as distributed processors~\cite{conway2008}.  Charged Conscious Points
(eCPs and qCPs) form stable clusters (``cages'') corresponding to SM particles.

This paper derives the SSV binding mechanisms, demonstrates geometric
stability for the principal cage types, and provides a worked example for
the electron.  It is the quantitative foundation for the mass generation
<<<<<<< HEAD
framework in Paper~2.
=======
framework in SM-2~\cite{abshier_sm2}.
>>>>>>> c085da0cc636ab5fe5c35ba45eac2ff84c41fcdc

\textbf{Note on the fourth cage (Version~6 correction):}
Previous versions of this paper mapped the top quark to a fullerene-like
cage of ${\sim}60$ vertices.  Exact computation of the 600-cell distance
<<<<<<< HEAD
shells (PS-1, 2026) shows that no 60-vertex shell exists.  The leading
candidate is the 30-vertex shell at $d^2 = 2$ (shell~3 of the 600-cell).
All tables and references to the fourth cage have been updated.
=======
shells (PS-1, 2026~\cite{ps1}) shows that no 60-vertex shell exists.  The
leading candidate is the 30-vertex shell at $d^2 = 2$ (shell~3 of the
600-cell).  All tables and references to the fourth cage have been updated.
>>>>>>> c085da0cc636ab5fe5c35ba45eac2ff84c41fcdc

%=====================================================================
\section{Dipole Sea as Organisational Medium}
%=====================================================================

The Dipole Sea consists of randomly oriented dipole pairs (DPs) that
oscillate at ZBW frequencies $f_{\text{ZBW}} \approx 1/(2 t_{\text{Pl}})$.
Mass and binding energy emerge as departures from this randomness: stable
particle structures condense local organisation, compressing the sea into
coherent patterns while globally minimising SSV disruptions.

%=====================================================================
\section{SSV Gradient Force Law}
%=====================================================================

The SSV at any Grid Point is the vector sum of contributions from all
CPs within its Planck Sphere Radius:

\begin{equation}
\vec{\text{SSV}}_{\text{GP}} =
\sum_{i \in \text{PSR}_i}
\frac{\text{SSV}_{0,i} \cdot p_i \cdot t_i \cdot
      f(\text{type compatibility}) \cdot \hat{r}_{i \to \text{GP}}}
     {r_{i \to \text{GP}}^2}
\label{eq:ssv_field}
\end{equation}

where $\text{SSV}_{0,i}$ is the elementary stress magnitude, $p_i = \pm 1$
is polarity, $t_i$ is the type coupling factor (eCP: $t=1$; qCP: $t=0.5$,
reflecting the colour-like rotational symmetry of the 600-cell --- this
<<<<<<< HEAD
factor is derived rigorously in the Strong Sector paper, Theorem~1), and
=======
factor is derived rigorously in SS-1~\cite{abshier_ss}, Theorem~1), and
>>>>>>> c085da0cc636ab5fe5c35ba45eac2ff84c41fcdc
$f(\text{type compatibility})$ modulates eCP--qCP interactions.

The force on a test CP $j$ is:
\begin{equation}
\vec{F}_j = -q_j \cdot \vec{\text{SSV}}(\vec{r}_j)
\end{equation}

Opposite polarities attract; like polarities repel.

%=====================================================================
\section{Potential Energy of Cage Configurations}
%=====================================================================

The total binding energy of a cage is:
\begin{equation}
E_{\text{binding}} = -\frac{1}{2} \sum_{j} q_j \cdot \Phi(\vec{r}_j)
\end{equation}

where $\Phi(\vec{r}_j) = \sum_{i \neq j}
\text{SSV}_{0,i} \cdot p_i \cdot t_i \cdot
f(\text{type}) / |\vec{r}_i - \vec{r}_j|$.

Stable cages minimise $E_{\text{binding}}$ (maximise its magnitude)
by aligning polarities and types along lattice symmetries that minimise
average distance and maximise constructive SSV interference.

%=====================================================================
\section{Geometric Preference for Cage Symmetries}
%=====================================================================

The 600-cell naturally supports the following cage subgraphs:

\begin{itemize}
\item \textbf{Tetrahedral cage} (4 vertices): 4 of the 12 nearest
  neighbours in tetrahedral arrangement.  Minimal stable shell.
\item \textbf{Icosahedral cage} (12 vertices): full first distance
  shell ($d^2 = 1/\varphi^2$).  Used for charm quarks, tau leptons, Z boson.
\item \textbf{Dodecahedral cage} (20 vertices): second distance shell
  ($d^2 = 1$).  Used for bottom quarks, Higgs boson.
\item \textbf{30-vertex shell} (30 vertices, $d^2 = 2$, shell~3):
<<<<<<< HEAD
  \textbf{replaces the C$_{60}$ assignment} (PS-1, 2026).  All 30 vertices
  equidistant from the apex, degree~4, vertex-transitive.  Leading candidate
  for the top quark fourth cage.  Mass formula using this shell is open
  (OP-SS-1).
\end{itemize}

\begin{figure}[h]
\centering
\caption{Tetrahedral cage: central eCP (electron) surrounded by 4
  compensating CPs at tetrahedral vertices (4 of 12 nearest 120CPs).}
\label{fig:tetrahedral-cage}
\end{figure}

=======
  \textbf{replaces the C$_{60}$ assignment} (PS-1, 2026~\cite{ps1}).
  All 30 vertices equidistant from the apex, degree~4, vertex-transitive.
  Leading candidate for the top quark fourth cage.
  Mass formula using this shell is open (OP-SS-1).
\end{itemize}

>>>>>>> c085da0cc636ab5fe5c35ba45eac2ff84c41fcdc
%=====================================================================
\section{Quantitative Cage Comparisons}
%=====================================================================

Binding energies scale with shell occupancy and geometric closure.
The values below are \emph{approximations} in lattice units, using the
relation $E_{\text{binding}} \approx N/2$ (from the worked electron
example below), not derived cage-by-cage.  The 30-vertex shell entry
(N=30) gives $E \approx 15$, replacing the previous C$_{60}$ (N=60)
entry of 30.

\begin{table}[h]
\centering
\begin{tabular}{lcccc}
\toprule
Cage Type & Vertices ($N$) & Binding ($\approx N/2$) & Example Particles
  & 600-cell shell \\
\midrule
Tetrahedral & 4 & 2.0 & $e$, $\mu$, $u$, $d$ & shell~0 subset \\
Icosahedral & 12 & 6.0 & charm, $\tau$, Z & shell~0 full \\
Dodecahedral & 20 & 10.0 & bottom, Higgs & shell~1 \\
30-vertex shell & 30 & ${\sim}15$ & top (candidate) & shell~3 \\
\midrule
<<<<<<< HEAD
\multicolumn{5}{l}{\small\textit{Note: C$_{60}$ (60 vertices) removed --- no such distance
  shell exists in the 600-cell (PS-1, 2026).}} \\
\bottomrule
\end{tabular}
\caption{Approximate binding energies for cage types (lattice units).
  Values use $E \approx N/2$; this is an approximation, not a per-cage derivation.
  The 30-vertex shell entry replaces the previously listed C$_{60}$ cage.}
=======
\multicolumn{5}{l}{\small\textit{Note: C$_{60}$ (60 vertices) removed ---
  no such distance shell exists in the 600-cell (PS-1, 2026~\cite{ps1}).}} \\
\bottomrule
\end{tabular}
\caption{Approximate binding energies for cage types (lattice units).
  Values use $E \approx N/2$; this is an approximation, not a per-cage
  derivation.  The 30-vertex shell entry replaces the previously listed
  C$_{60}$ cage.}
>>>>>>> c085da0cc636ab5fe5c35ba45eac2ff84c41fcdc
\label{tab:cage_binding}
\end{table}

%=====================================================================
\section{Stability Conditions}
%=====================================================================

Partial occupancy leads to instability:

\begin{table}[h]
\centering
\begin{tabular}{ccc}
\toprule
Shell Occupancy & Stability & Reason \\
\midrule
1 CP & Unstable & Strong residual gradient \\
2 CPs & Metastable & Residual torque remains \\
3 CPs & Highly unstable & Asymmetric gradients dominate \\
4 CPs & Minimal stable & Tetrahedral symmetry cancels gradients \\
\bottomrule
\end{tabular}
\caption{Stability as a function of tetrahedral shell occupancy.}
\label{tab:stability}
\end{table}

%=====================================================================
\section{Worked Example: Electron Binding Energy}
%=====================================================================

\subsection{Setup}

Electron model: central $-$eCP, surrounded by 4 positive compensating
CPs in the tetrahedral shell.  Parameters: $\text{SSV}_0 = 1$ (natural
lattice units), nearest-neighbour distance $d_0 = 1$, eCP--eCP coupling
factor $t = 1$.

\subsection{Potential and Binding Energy}

Potential at the central CP from one shell CP: $\Phi = -\text{SSV}_0/d_0 = -1$.
For four tetrahedral CPs: $\Phi_{\text{total}} = -4$.

Binding energy:
\begin{equation}
E_{\text{binding}}
= -\tfrac{1}{2} \times (-4) \times \text{SSV}_0
= 2 \quad (\text{lattice units})
\label{eq:electron_binding}
\end{equation}

\subsection{Calibration to Physical Units}

\textbf{This step is a calibration, not a derivation.}
The binding energy of 2 lattice units was constructed to reproduce the
electron rest mass.  Setting $E_{\text{binding}} \times \text{SSV}_0 = m_e c^2$
gives:

\begin{equation}
\text{SSV}_0
= \frac{m_e c^2}{2}
= \frac{0.511 \, \text{MeV}}{2}
= 0.2555 \, \text{MeV}
\label{eq:ssv0_value}
\end{equation}

This is the one calibration constant of the binding-energy framework.
All other particles use this same $\text{SSV}_0$, so the hierarchy is
determined by cage geometry alone --- but the absolute scale is set by
the electron.

<<<<<<< HEAD
Note: the K3 Spectral Theorem (Paper~3) provides a complementary derivation
of the lepton mass scale from the ZBW thermal energy $\hbar\omega_0 =
\text{sea\_strength} \times \hbar c / r_{\rm conf}$, where
$\text{sea\_strength} \approx 0.1780$ is derived from 600-cell geometry.
The relationship between $\text{SSV}_0 = 0.2555$~MeV and the thermal
energy scale is an open connection (OP-SS-1).
=======
The K3 Spectral Theorem (SM-3~\cite{abshier_paper3}) provides a
complementary derivation of the lepton mass scale from the ZBW thermal
energy $\hbar\omega_0 = \text{sea\_strength} \times \hbar c / r_{\rm conf}$,
where $\text{sea\_strength} \approx 0.1780$ is derived from 600-cell
geometry~\cite{abshier_ss}.  The relationship between
$\text{SSV}_0 = 0.2555$~MeV and the thermal energy scale is an open
connection (OP-SS-1).
>>>>>>> c085da0cc636ab5fe5c35ba45eac2ff84c41fcdc

%=====================================================================
\section{Connection to ZBW Spectrum and Geometric Suppression}
%=====================================================================

<<<<<<< HEAD
This paper's binding energies preview the unified ZBW spectrum in Paper~2.
For bound cages ($d=0$), full lattice coupling applies.  Linear ZBW extras
($d=1$) and unbound modes ($d=3$) receive suppression $\sigma = 120^{-d}$,
explaining neutrino masses.
=======
This paper's binding energies preview the unified ZBW spectrum in
SM-2~\cite{abshier_sm2}.  For bound cages ($d=0$), full lattice coupling
applies.  Linear ZBW extras ($d=1$) and unbound modes ($d=3$) receive
suppression $\sigma = 120^{-d}$, explaining neutrino masses.
>>>>>>> c085da0cc636ab5fe5c35ba45eac2ff84c41fcdc

The ZBW oscillation, reinterpreted from Dirac~\cite{dirac1928} as the
fundamental oscillatory mode of DPs in the Dipole Sea at frequency
$f_{\text{ZBW}} \approx 1/(2 t_{\text{Pl}})$, provides the microscopic
basis for mass as condensed organisational energy.

Forward references: the coupling factor $t = 0.5$ for qCPs is derived
<<<<<<< HEAD
rigorously in the Strong Sector paper (Theorem~1: SU(3) from tetrahedral
hopping operators).  The log-hierarchy formula in Paper~1c uses
$\varphi$-scaling motivated by this paper.
=======
rigorously in SS-1~\cite{abshier_ss} (Theorem~1: SU(3) from tetrahedral
hopping operators).  The log-hierarchy formula in SM-TN-2~\cite{abshier_tn2}
uses $\varphi$-scaling motivated by this paper.
>>>>>>> c085da0cc636ab5fe5c35ba45eac2ff84c41fcdc

%=====================================================================
\section{Conclusion}
%=====================================================================

This paper derives the SSV binding mechanism, demonstrates that specific
cage geometries (tetrahedral, icosahedral, dodecahedral, 30-vertex shell)
are energetically stable, and provides a quantitative worked example for
the electron.

\textbf{One calibration constant} ($\text{SSV}_0 = 0.2555$~MeV, fixed to
the electron mass) sets the absolute scale.  The mass hierarchy is then
determined geometrically by cage vertex count and shell structure.

The fourth cage for the top quark is now identified as the 30-vertex
<<<<<<< HEAD
shell at $d^2 = 2$ of the 600-cell (PS-1, 2026), replacing the
=======
shell at $d^2 = 2$ of the 600-cell (PS-1, 2026~\cite{ps1}), replacing the
>>>>>>> c085da0cc636ab5fe5c35ba45eac2ff84c41fcdc
falsified C$_{60}$ assignment.  The mass formula using this cage
remains an open problem (OP-SS-1).

\appendix

%=====================================================================
\section{Holographic Constraints, DI-Bit Foundation, and Lattice Addressability}
%=====================================================================

The total number of Grid Points within any 120CP vertex is constrained
by the holographic principle.  A rough estimate is ${\sim}10^{60}$
addressable positions per vertex (motivated by the Bekenstein bound
for a Planck-sphere volume), giving ${\sim}10^{61}$ total across all 120
vertices within the cosmic horizon.  \textbf{This is an order-of-magnitude
estimate, not a derived result.}

The SSV field emerges from DI-bit exchange between 120CPs:
\begin{equation}
\text{SSV}_{0,i} = \alpha \cdot \frac{\Delta b_i}{t_{\text{Pl}}}
\end{equation}
where $\Delta b_i$ is the net DI-bit excess at 120CP $i$ and $\alpha$
is the bit-to-stress conversion factor.

%=====================================================================
\section{Quantitative Extensions, Failure Modes, and Speculative Signatures}
%=====================================================================

\subsection{Binding Energy Table (updated)}

See Table~\ref{tab:cage_binding} in the main text, which now lists the
30-vertex shell (N=30, $E \approx 15$) as the fourth cage, replacing
the C$_{60}$ entry.

\subsection{Failure Modes and Instabilities}

Partial occupancy or type mismatch leads to instabilities:
\begin{itemize}
\item Partial tetrahedral occupancy (1--3 CPs): dissociation or reconfiguration.
\item Resonance: DP oscillation frequencies aligned with lattice modes
  $\to$ destabilisation or energy radiation.
\item Type mismatch: excessive qCP--qCP interactions ($f=0.5$) can collapse
  into glueball-like states.
\end{itemize}

\subsection{Lattice Discretisation Effects}

Angular quantisation restricts orientations to icosahedral symmetries.
Continuous limits emerge statistically from large $N_k$ ensembles.
Finite granularity sets a natural ultraviolet cutoff.

\subsection{Speculative Experimental Signatures}

\textbf{The following are speculative estimates without full derivations.}
They are registered as open problems rather than predictions:
\begin{itemize}
\item Ultra-precise electron $g$-2 measurements might show deviations of
  ${\sim}10^{-12}$ from QED due to finite lattice effects (not derived;
  see OP-QM series).
\item High-energy resonances at ${\sim}10^{16}$~eV spaced by
  $\Delta E \sim \text{SSV}_0 \times \varphi^n$ (open problem).
\item Vacuum birefringence at ${\sim}10^{-15}$ in the Cotton-Mouton
  effect (not derived; speculative).
\end{itemize}

These require full derivation before they can be considered predictions.

%=====================================================================
\section*{Acknowledgements}
%=====================================================================

Thomas Lee Abshier ND conceived the CPP framework and directed the
research programme.  Grok (xAI) provided computational collaboration,
theoretical refinement, and simulation support.  Claude Sonnet (Anthropic)
performed the technical review that produced Version~6, identifying the
C$_{60}$ correction, the SSV$_0$ calibration clarification, and the other
corrections noted in the abstract.

%=====================================================================
\begin{thebibliography}{9}

\bibitem{pdg2024}
Particle Data Group,
Review of Particle Physics,
\textit{Prog.\ Theor.\ Exp.\ Phys.} \textbf{2024}, 083C01 (2024).

\bibitem{gellmann1964}
M.~Gell-Mann,
A Schematic Model of Baryons and Mesons,
\textit{Phys.\ Lett.} \textbf{8}, 214 (1964).

\bibitem{zweig1964}
G.~Zweig,
An SU(3) Model for Strong Interaction Symmetry and its Breaking,
CERN-TH-401 (1964).

\bibitem{dirac1928}
P.~A.~M.~Dirac,
The Quantum Theory of the Electron,
\textit{Proc.\ Roy.\ Soc.\ Lond.\ A} \textbf{117}, 610 (1928).

\bibitem{conway2008}
J.~H.~Conway and N.~J.~A.~Sloane,
\textit{Sphere Packings, Lattices and Groups}, 3rd ed.,
Springer (2008).

\bibitem{abshier_ss}
T.~L.~Abshier, Grok (xAI), Claude Sonnet (Anthropic),
<<<<<<< HEAD
Conscious Point Physics: The Strong Sector from the 600-Cell Lattice
(Strong Sector paper, Version~2),
CPP Series, 2026.

\bibitem{abshier_paper3}
T.~L.~Abshier, Grok (xAI), Claude Sonnet (Anthropic),
K3 Spectral Theorem and the Koide Formula (Paper~3),
CPP Series, 2026.
=======
Conscious Point Physics: The Strong Sector from the 600-Cell Lattice (SS-1),
\textit{600-Cell Standard Model Emergence Series}, 2026.\\
\url{https://github.com/Hyperphysics-Institute/CPP/blob/main/series_strong/SS-1_strong_sector_from_600cell_lattice.tex}

\bibitem{abshier_sm2}
T.~L.~Abshier, Grok (xAI), Claude Sonnet (Anthropic),
Mass Generation from Geometric Hierarchies and Cage Complexity (SM-2),
\textit{600-Cell Standard Model Emergence Series}, 2026.\\
\url{https://github.com/Hyperphysics-Institute/CPP/blob/main/series_standard_model/papers/SM-2_mass_generation_geometric_hierarchies.tex}

\bibitem{abshier_paper3}
T.~L.~Abshier, Grok (xAI), Claude Sonnet (Anthropic),
K3 Spectral Theorem and the Koide Formula (SM-3),
\textit{600-Cell Standard Model Emergence Series}, 2026.\\
\url{https://github.com/Hyperphysics-Institute/CPP/blob/main/series_standard_model/papers/SM-3_k3_spectral_theorem_koide_formula.tex}

\bibitem{abshier_tn2}
T.~L.~Abshier, Grok (xAI), Claude Sonnet (Anthropic),
Reconstruction 2: Bridge from Original to 600-Cell Framework (SM-TN-2),
\textit{600-Cell Standard Model Emergence Series}, 2026.\\
\url{https://github.com/Hyperphysics-Institute/CPP/blob/main/series_standard_model/papers/SM-TN-2_bridge_original_to_600cell.tex}

\bibitem{ps1}
T.~L.~Abshier, Grok (xAI), Claude Sonnet (Anthropic),
600-Cell Distance Shell Analysis (PS-1): Exact shell volumes and
falsification of the C$_{60}$ cage assignment,
\textit{600-Cell Standard Model Emergence Series}, 2026.\\
\url{https://github.com/Hyperphysics-Institute/CPP/blob/main/series_standard_model/notebooks/ps1_quark_mass_ladder.ipynb}
>>>>>>> c085da0cc636ab5fe5c35ba45eac2ff84c41fcdc

\end{thebibliography}

\end{document}
